% Proposed B.114. Isolated owned component; no assembled/register edit.
\subsection*{The actual annular neck and its two additional relations}

\statement{B.114}{Proved}{The actual annular one-neck group is a proper-amalgam quotient by two nontrivial based cross commutators}
\label{stmt:B114}
Use the actual surface, involution and first-neck marking of B.106,
B.108 and B.112. Retain the complete annuli
\[
 B_-=\{0<|q|<1\},\qquad B_+=\{1<|q|<\infty\},
\]
including every inner and outer critical fiber. Join them through
one sufficiently small regular neck at $q_*=\exp(\pi i/5)$, with
contractible overlap, and denote the resulting base by $U_{\rm ann}$.
No other unit-circle neck is included. Let $R_-$ and $R_+$ be the
actual complete bad-region orbifold groups on the inner and outer
sides. Their cones are $c_-$ at $t=\infty$ and $c_+$ at $t=0$,
with the common neck whiskers fixed below. Let $H_\pm$ be their
ordinary index-two subgroups. The actual common disk subgroup
$K\cong F_5$ embeds in both $H_\pm$ and is a retract of each.

Put $\mathcal A=R_-*_{K}R_+$. This abstract amalgam and its full
profinite completion satisfy
\begin{equation}\label{eq:B114-proper}
 \widehat{\mathcal A}\cong
 \widehat R_-\amalg_{\widehat K}\widehat R_+,
\end{equation}
with injective completed vertex and edge maps. If $L_\pm\cong F_5$
are the actual bad fixed-source images of B.112, transported with
the same cone whiskers, then
\begin{equation}\label{eq:B114-abstract-centralizers}
 \begin{aligned}
 C_{\mathcal A}(c_\pm)&=C_{R_\pm}(c_\pm)
                         =L_\pm\times\langle c_\pm\rangle,\\
 C_{\widehat{\mathcal A}}(c_\pm)
     &=C_{\widehat R_\pm}(c_\pm)
                         =\widehat L_\pm\times\langle c_\pm\rangle.
 \end{aligned}
\end{equation}

Let $k_\pm$ be the actual boundary-section loops on the two annuli.
Their geometric base winding is one. The actual orbifold group
$G_{\rm ann}$ of the entire one-neck region is instead
\begin{equation}\label{eq:B114-actual-quotient}
 \begin{split}
 G_{\rm ann}&\cong
 \mathcal A/\langle\!\langle
       [c_+,k_-],\ [c_-,k_+]\rangle\!\rangle,\\
 \widehat G_{\rm ann}&\cong
 \widehat{\mathcal A}/
 \overline{\langle\!\langle[c_+,k_-],\ [c_-,k_+]\rangle\!\rangle}.
 \end{split}
\end{equation}
The bar means closed normal closure in the full profinite group.
Each of these two relators is nontrivial in $\mathcal A$ and in
$\widehat{\mathcal A}$. Thus neither arrow in
\eqref{eq:B114-actual-quotient} is injective. More strongly, either
relator remains nontrivial after killing the other, both discretely
and after taking its closed normal closure in the full completion.

The selected quotient source over $U_{\rm ann}$ has abstract
fundamental group $F_6$; its actual connected original source cover
has group $F_{11}$. The selected inner source loop is $k_-$, and
its image in $G_{\rm ann}$, or in $\widehat G_{\rm ann}$, does not
belong to the corresponding image of the outer bad vertex.
These source counts do not assert ambient source injectivity.
Nor do we assert injectivity of a bad vertex after the two
relations, or a centralizer formula for the final quotient.

\begin{proof}
\emph{B.114.1. One actual neck basis and its annular retractions.}
Use the lower disk containing the four literal $A$ punctures and
the selected branch puncture in B.108. Its disk, five disjoint
meridian arcs and single gap connector are transported together
through the first regular neck. At each endpoint the disk is
isotoped within the lower half-plane into the corresponding bad
side, fixing the gap. The five punctures in that half-plane are
exactly the selected ones, so this is a based disk comparison,
not a list of independently chosen peripheral conjugacy classes.

Here is the explicit comparison with B.112's actual inner marking.
Write its adapted variables as $(a,b,C,f,d,e)$ and put
$\bar x=CxC$. We use $C$, rather than $c_-$, to distinguish this
adapted cone from the already fixed neck cone. In the original
six-strand marking the common disk meridians are
$x_1,x_3,x_4,x_5,x_6$. Their literal expressions and inverses are
\begin{equation}\label{eq:B114-neck-basis}
 \begin{aligned}
 x_1&=a\bar b a^{-1},&x_3&=a,&x_4&=fde d^{-1}f^{-1},\\
 x_5&=d,&x_6&=d^{-1}fd,&&\\
 a&=x_3,&\bar b&=x_3^{-1}x_1x_3,&d&=x_5,\\
 f&=x_5x_6x_5^{-1},&
 e&=x_6^{-1}x_5^{-1}x_4x_5x_6.&&
 \end{aligned}
\end{equation}
Both compositions are free-group identities. Therefore this
actual disk is generated equally by $(a,\bar b,d,e,f)$.
The complete annular ordinary detector of B.112 sends this tuple
to the free basis $(X_1,X_2,X_3,X_4,H)$ of $F_5$. It follows that
the disk map is injective. Identifying the detector target with
this actual subgroup yields a retraction
\[
 r_-:H_-\longrightarrow K,\qquad r_-|_K=\operatorname{id}.
\]
The detector was checked against all ordinary node relations and
both directions of the retained zero-loop action. This is a
retraction of the complete annular group, not only its generic
fiber or slit quotient.

The computational base is connected to the first neck by the
regular path used in B.108 and B.110. Along its radial part the
four selected moving punctures and the selected branch stay in
the same half-plane; no selected pair collides. Along the first
unit-circle arc they remain distinct in the lower semicircle.
Transporting the whole disk and connector gives exactly the same
based subgroup. At $q_*$ choose its annular-sector representative
with logarithmic radial thickness symmetric about the unit circle
and its based gap on that circle. It is invariant under
\[
 J(q,t)=(1/\bar q,1/\bar t).
\]
The actual identity $T_u(1/\bar q)=1/\overline{T_u(q)}$ preserves
the physical graphs. At $q_*$ the selected five points lie on
the unit circle and are individually fixed by this reflection.
Let $\eta=J_*|_K$ be its based disk automorphism. Transporting the
outer marking and connector by $J$ gives
\[
 r_+=\eta^{-1}r_-J_*:H_+\longrightarrow K,
 \qquad r_+|_K=\operatorname{id}.
\]
This does not assume that $\eta$ inverts five free generators
separately; it retains the actual reflected whiskers.

For later precision, the inner detector satisfies
\begin{equation}\label{eq:B114-nonzero-nu}
 r_-\operatorname{Inn}(C)|_K:
 (a,\bar b,d,e,f)\longmapsto(1,1,1,1,f).
\end{equation}
It is not the zero homomorphism. The old comb cone is
$c_{\rm old}=P C P^{-1}$ with $P=a\bar b^{-1}\in K$.
A neck cone whisker may involve a further common based change.
We retain that change throughout rather than replace only its
cone letter by $C$. The following argument works for the actual
cone/retraction endomorphism, whatever these retained prefixes.

\emph{B.114.2. Fully invariant finite edge kernels.}
We prove the abstract comparison needed here. Let a finitely
generated group $K$ be a retract of each ordinary subgroup
$H_i$ in $R_i=H_i\rtimes\langle c_i\rangle$, with $c_i^2=1$.
Write $r_i:H_i\to K$ for the retractions and
\[
 \nu_i=r_i\operatorname{Inn}(c_i)|_K\in\operatorname{End}(K).
\]
Neither triviality nor involutivity of $\nu_i$ is required.
For $N\triangleleft K$ of finite index, put $Q=K/N$ and
\begin{equation}\label{eq:B114-fully-invariant}
 L=\bigcap_{\phi\in\operatorname{Hom}(K,Q)}\ker\phi.
\end{equation}
There are finitely many such homomorphisms since $K$ is finitely
generated. Hence $L$ is normal of finite index; the quotient map
is among them, so $L\subset N$. For every endomorphism $\theta$
the homomorphisms $\phi\theta$ occur in the same set. Thus
$\theta(L)\subset L$: $L$ is fully invariant.

Put $Q_L=K/L$. The homomorphism on $H_i$ defined by
\[
 x\longmapsto(r_i(x)\bmod L,\ r_i(c_i x c_i)\bmod L)
\]
is equivariant for coordinate interchange. It extends to
\begin{equation}\label{eq:B114-paired}
 \Phi_i:R_i\longrightarrow
 (Q_L\times Q_L)\rtimes_{\rm swap}C_2,
 \qquad c_i\longmapsto\text{the interchange involution}.
\end{equation}
On $K$ it is $k\mapsto(k\bmod L,\nu_i(k)\bmod L)$.
Its kernel on $K$ is exactly $L$: the first coordinate forces
membership in $L$, and full invariance kills the second there.
This is why no vanishing assertion about $r_i(c_i Kc_i)$ is needed.

Given arbitrary finite-index normal subgroups $M_i$ of the two
vertices, choose $N\subset(M_-\cap K)\cap(M_+\cap K)$, construct
$L$ by \eqref{eq:B114-fully-invariant}, and set
$N_i=M_i\cap\ker\Phi_i$. Then
\[
 N_i\triangleleft R_i,
 \qquad [R_i:N_i]<\infty,
 \qquad N_i\subset H_i,
 \qquad N_-\cap K=N_+\cap K=L.
\]
The last equality follows also from $L\subset M_i$.
Thus arbitrary vertex quotients admit compatible refinements with
the same exact edge kernel. The two finite edge embeddings are
their actual graph maps; they need not be the same map to a
product in \eqref{eq:B114-paired}.

\emph{B.114.3. Proper full completion of the abstract amalgam.}
Apply the preceding construction to the actual retractions. Put
$A_i=R_i/N_i$, $D=K/L$, and $V=A_-*_{D}A_+$.
This finite-group amalgam is residually finite and virtually free
of finite rank. For completeness, choose a finite cardinal
divisible by both $|A_i|$, and realize both vertices by repeated
regular actions on sets of that cardinal. Their restrictions are
isomorphic free $D$-sets, so identify them. The resulting finite
permutation representation of $V$ is injective on both vertices.
Its kernel acts freely on the ordinary amalgam tree and is free.
It has finite rank because it has finite index in the finitely
generated group $V$. Residual finiteness of this free kernel and
finite-index cores then prove residual finiteness of $V$.

Every prescribed finite quotient of either original vertex is
retained after refinement in a finite quotient of $\mathcal A$.
Consequently both vertex completions inject in
$\widehat{\mathcal A}$. Each edge completion already embeds in
its ordinary vertex by retraction, and in the orbifold vertex by
index-two cores. The universal property on finite targets gives
\eqref{eq:B114-proper}, with properness proved. Also every finite
quotient of $\mathcal A$ factors through one of the $V$: start
with the kernels of its two vertex restrictions. These arguments
do not require residual finiteness of either original vertex.

\emph{B.114.4. The odd-cone centralizer before the extra relations.}
In the ordinary amalgam tree a selected odd cone cannot lie in a
conjugate edge group, since the edge is even. It fixes just its
base vertex. Hence
$C_{\mathcal A}(c_i)=C_{R_i}(c_i)$.
The same argument in each finite-group amalgam $V$ gives
\[
 C_V(c_i)=C_{A_i}(c_i),
\]
which is a finite subgroup.

We verify the completed step rather than assume a profinite tree
assertion for the unknown vertices. A normal finite-index free
subgroup $F_r$ of $V$ has its full topology by finite-index cores.
The abstract and completed extensions have the same finite
quotient and kernels $F_r,\widehat F_r$. For every discrete
continuous $\mathbb F_p[[\widehat F_r]]$-module, degree-one
cohomology agrees: a cocycle is specified on a finite free basis;
its finitely many values and their compact orbits span a finite
invariant submodule, and the associated finite semidirect product
extends the cocycle continuously. Coboundaries agree, degree zero
agrees by density, and higher cohomology vanishes on both sides.
The two Lyndon--Hochschild--Serre spectral sequences therefore
give the comparison for $V$ with every such discrete module,
with only free-subgroup rows zero and one. For finite modules
both rows, and their finite-quotient cohomology in every degree,
are finite. Thus $V$ is residually finite, has finite virtual
$p$-cohomological \emph{type}, and has the required all-module
cohomology comparison. The complete prime-order theorem of
\citet[Theorem~4.3, Sections~4.1--4.2]{BoggiZalesskii2025},
with the hypotheses checked as in B.108, applies at $p=2$.
It gives
\[
 C_{\widehat V}(c_i)=\overline{C_V(c_i)}=C_{A_i}(c_i).
\]

If $x\in\widehat{\mathcal A}$ centralizes $c_i$, its image in
every completed head belongs to the corresponding finite vertex.
If $x$ lay outside the compact image of $\widehat R_i$, a finite
quotient would separate it from that image. Factoring this quotient
through one of the heads contradicts the preceding conclusion.
Therefore $C_{\widehat{\mathcal A}}(c_i)=C_{\widehat R_i}(c_i)$.
The actual source and topology statement B.112 now gives
\eqref{eq:B114-abstract-centralizers} for the transported cones.
This argument concerns $\mathcal A$, before any cross commutator
is killed.

\emph{B.114.5. The actual categorical pushout and its collar loops.}
The common orbifold neck fiber, with the one transported disk
and its two cone whiskers, has group
\[
 H_{\rm neck}=C_2(c_+)*C_2(c_-)*K.
\]
Over a complete annulus the good-side contraction of B.107 gives
the product of its cone with the base. Its intersection with the
bad side is the actual boundary collar, with infinite cyclic
fundamental group. Write its based generator in the bad vertex
as $k_i$. The geometric $q$-projection sends $k_i$ to a generator
of the annular base and retracts onto this cyclic subgroup. The
two entire regional orbifold groups are consequently
\begin{equation}\label{eq:B114-regional}
 \begin{aligned}
 E_-&=(C_2(c_+)\times\langle k_-\rangle)
                      *_{\langle k_-\rangle} R_-,\\
 E_+&=(C_2(c_-)\times\langle k_+\rangle)
                      *_{\langle k_+\rangle} R_+.
 \end{aligned}
\end{equation}
The good-side base circle cannot be removed from this calculation.

We identify its section word explicitly. In the inner bad quotient
use $y=z^2$, $z=1/t$. The finite moving punctures
$0,q,h_j(q)^2$ lie in $|y|<1$; the retained boundary point is
$y=1$. The plane-to-disk trivialization used for the actual far
braids sends the far section to a point at constant argument
$\theta\in(-\pi,0)$ and radius greater than
\[
 \rho(q)=\frac{1+\max(|q|,|h_j(q)^2|)}2.
\]
This radius is strictly greater than the modulus of every finite
puncture, for every $q$ in the annulus, including its retained
critical fibers. First move within the empty outer annulus to a
continuous radius $r(q)$ with $\rho(q)<r(q)<1$. Rotate from
argument $\theta$ to $+\pi$ at that radius, and only then move
radially out to $y=-1$. In particular the positive real ray is
crossed strictly inside the unit circle; the retained point
$y=1$ is avoided. This is a homotopy of sections over the whole
base. With the specified principal lift it ends at $z=+i$ and
$t=-i$, the selected lower-$t$ gap. A rotation to $-\pi$ would
instead end at $t=+i$ and cannot be substituted here.
The endpoint is the quotient of the two exchanged boundary
corridors. Whisker the good cone through its empty product disk
from this endpoint. The endpoint base loop commutes with that
cone. The section homotopy identifies this loop with the original
far loop by its single common connector. Thus
\[
 k_-=g_0=f^{-1}h
\]
in B.112's consistently transported adapted marking. Its geometric
base winding is one, whereas B.112's unfolding height sends this
same word to zero. Reflect this complete section homotopy and
connector by $J$ to define $k_+$ on the other side. Reversing the
chosen orientation of $k_+$ would not change its commuting relation.

We also give the concrete cone word for this selected connector.
At the actual computational base $q_b=4/5+i/2$, use the original
cut coordinate $X=(1+371i/1000)y$, downward cuts, and far base
$X=10$. Put $\epsilon=1/1000$. A polygon for the upper cone
whisker has consecutive vertices
\[
 (10,0),\ (10,5),\ (-2,5),\ (-2,-742/1000),\
 (-\epsilon,-371\epsilon/1000),\ (-\epsilon,0).
\]
The exact ordered cut word of this whisker is $x_1$; the positive
small square around the cone has word $x_2$. Every segment and
crossing has a strict rational separation certificate, preserved
by the previously certified actual-root displacement. The full
based cone is therefore
\begin{equation}\label{eq:B114-actual-gap-cones}
 c_-=x_1x_2x_1^{-1}=aC_-a^{-1},\qquad
 c_+=\eta^{-1}(a)C_+\eta^{-1}(a)^{-1}.
\end{equation}
The first equality in adapted variables reduces literally using
B.112's inverse basis and $C_-^2=1$ only. Here $C_-=C$ and $C_+$
is its reflected adapted cone; both words are expressed through
the common $K$ embeddings. The direction $\eta^{-1}$ follows
from $J_*:R_+\to R_-$ and $J_*|_K=\eta$.

For the physical interpretation of the unbounded polygon, apply
the actual radial plane-to-disk homeomorphism first: its upper
outer portion lies inside $|y|<1$, above every moving radius,
and the retained boundary point stays outside. The five ordinary
punctures and their original common-entry comb system lie in
the upper $y$ half-plane along the chosen regular parameter
bridge. The upper connector transports that whole system to the
upper-half-plane disk group based at $y=-1$. At the unit neck,
$J(y)=1/\bar y$ preserves this half-plane and fixes the basepoint.
Thus this change preserves the actual selected $K$ and its
principal lift, not merely the cone conjugacy class. No formula
for $\eta$ on individual generators is assumed.

Orbifold van Kampen glues \eqref{eq:B114-regional} by the actual
two maps from $H_{\rm neck}$. First identify the two copies of
$K$ through the transported disk basis. Each free good cone is
then identified with its bad cone in the opposite vertex.
Eliminating these two redundant generators leaves precisely the
first presentation in \eqref{eq:B114-actual-quotient}.
The only new regional relations are the two displayed cross
commutators. This is categorical elimination; it does not assume
that the entire neck fiber injects after critical filling.
The bad-vertex amalgam preceding the quotient is legitimate by
Section~1. Completion gives the second presentation because its
continuous finite targets are exactly the finite targets of the
first presentation. No injection into this quotient follows.

The notation $c_\pm$ in this pushout always means the common-neck
cone with its actual whisker, concretely
\eqref{eq:B114-actual-gap-cones} in this gauge. It is different
from the old comb cone $P C P^{-1}$ with $P=a\bar b^{-1}$.
Replacing only a cone word by $C$ would change the assertion.
A common conjugation of an entire regional marking must also
conjugate its $K$ embedding and its section word. The categorical
presentation keeps these maps together.

\emph{B.114.6. Both new relators are genuinely killed elements.}
The base winding on $R_-$ is zero on $K$ and one on $k_-$.
Hence $k_-\notin K$. The element $c_+$ is odd, so $c_+\notin K$.
Therefore $[c_+,k_-]$ is a reduced alternating word of length
four in $\mathcal A$ and is nontrivial. The reflected argument
applies to $[c_-,k_+]$.

There are particularly small independent finite detectors. Map
$R_-$ by its geometric base winding to $\mathbb Z=\langle t\rangle$
and $R_+$ by orientation to $C_2=\langle c\rangle$.
Both maps kill $K$, so they induce
\[
 D_-:\mathcal A\longrightarrow\mathbb Z*C_2,
 \qquad (k_-,c_+,c_-,k_+)\longmapsto(t,c,1,1).
\]
The reflected assignment gives $D_+$ with
$(k_+,c_-,c_+,k_-)\mapsto(t,c,1,1)$.
Now compose each map with
\[
 \mathbb Z*C_2\longrightarrow S_3,
 \qquad t\longmapsto(123),\qquad c\longmapsto(12).
\]
With functional permutation composition, $ctc=t^{-1}$, and hence
$[c,t]=t\ne1$. Writing $r_-=[c_+,k_-]$ and
$r_+=[c_-,k_+]$, the two finite images are exactly
\begin{equation}\label{eq:B114-two-detectors}
 \begin{array}{c|cc}
       &r_-&r_+\\\hline
 D_-   &t&1\\
 D_+   &1&t.
 \end{array}
\end{equation}
Thus neither relator is in the normal closure of the other.
These finite maps extend to $\widehat{\mathcal A}$ and kill the
other closed normal closure, proving the same assertion there.
The exact permutation products, including the stated convention,
were checked directly in two prescribed maps; no finite-group
census is involved. No residual finiteness of an original vertex
is required for these witnesses.

\emph{B.114.7. The actual source and the extra base direction.}
The union $U_{\rm ann}$ consists of two annuli joined across a
contractible overlap. Its group is
\[
 \pi_1(U_{\rm ann})=F(a_0,b_0),
\]
where $a_0$ is the inner base loop and $b_0$ the outer one.
The selected quotient source is this base with the four outer
contact points removed. Its pieces are an inner annulus and an
outer annulus with four punctures, with the same contractible
overlap. Thus its abstract group has the geometric free basis
\[
 a_0,b_0,p_1,p_2,p_3,p_4.
\]
The four $p_j$ are the actual outer contact lassos. The inner
source loop maps to $k_-$, by the good product and the section
homotopy above. The other five map to the actual outer bad
fixed source. All six images centralize $c_+$ and are even.

The ambient $q$-projection induces
$G_{\rm ann}\to F(a_0,b_0)$: the neck group is fiberwise,
and each regional group maps to its corresponding annular
cyclic factor. In particular $k_-$ maps to $a_0$, while the
image of $R_+$ lies in $\langle b_0\rangle$. Hence the image
of $k_-$ cannot belong to that of $R_+$. The same strictness
holds after full completion: the finite base character
$a_0\mapsto1$, $b_0\mapsto0$ to $C_2$ separates $k_-$ from
the entire compact image of $\widehat R_+$. This proves the
extra source direction without a rank-only inference.

The original selected source cover adjoins $\sqrt{1-q}$. Its
character is zero on $a_0,p_1,\ldots,p_4$ and one on $b_0$:
the outer circle surrounds $q=1$, whereas the inner circle and
the small contact circles do not. It is consequently connected.
With coset representatives $1,b_0$, its free Schreier basis is
\[
 a_0,p_1,p_2,p_3,p_4,\quad b_0^2,\quad
 b_0a_0b_0^{-1},\quad b_0p_jb_0^{-1}\ (1\leq j\leq4).
\]
Its rank is $2(6-1)+1=11$. These statements concern the actual
source and its cover; they do not prove injectivity of their
ambient images. The effect of the two relations on the whole
centralizer remains a separate question, as does the addition
of the other four necks. Neither main decision is concluded.

\emph{B.114.8. Exact records.}
\begin{raggedright}
The literal neck basis, its inverse, cone conversion and actual
retraction are recorded in
\path{work_DE/next_marginal_source/annular_neck_retraction23.json},
with its producer and genuine output alongside. The complete
geometric/source proof is
\path{work_DE/next_marginal_source/annular_neck_retraction_and_source23.md}.
The based categorical pushout, section homotopy and cone safeguards
are in \path{work_CE/annular_first_neck_pushout23.txt}.
The corrected selected-sheet cut and basis certificate is
\path{work_CE/next_symmetric_nodes/annular_upper_gap_cone23.json},
with producer and genuine output
\path{work_CE/code/check_annular_upper_gap_cone23.py} alongside.
The earlier lower-corridor computation is preserved as a
differently lifted connector; it is not used for the selected gap.
The independent abstract finite-kernel proof is
\path{work_A/fully_invariant_paired_kernel_lemma23.md}.
The two exact permutation witnesses are in
\path{work_A/B114_cross_relators23.json}, with producer
\path{work_A/code/check_B114_cross_relators23.py} and genuine output
alongside. The earlier valid finite-head argument for detecting the
two relators is preserved in
\path{work_A/B114_finite_head_cross_proof23.tex}.
No numerical group enumeration is used in this statement.\qedhere
\par\end{raggedright}
\end{proof}
